\documentclass[11pt,a4paper]{ctexart}
\usepackage[margin=2.4cm]{geometry}
\usepackage{amsmath,amssymb}
\usepackage{fancyhdr}
\usepackage{enumitem}
\linespread{1.18}
\setlength{\headheight}{15pt}
\setlength{\parskip}{0pt}

\pagestyle{fancy}
\fancyhf{}
\fancyhead[L]{\small Team 2300336}
\fancyhead[R]{\small \thepage}
\renewcommand{\headrulewidth}{0.4pt}
\fancypagestyle{plain}{\fancyhf{}\renewcommand{\headrulewidth}{0pt}}

\ctexset{
  section/format = \Large\bfseries\raggedright,
  subsection/format = \large\bfseries\raggedright,
  subsubsection/format = \normalsize\bfseries\raggedright
}
\setcounter{secnumdepth}{3}

\begin{document}

%====================== 摘要页 ======================
\thispagestyle{empty}
\begin{center}
{\small 2023 MCM/ICM\quad 选题 A\quad 队伍控制号 2300336\quad 摘要页（Summary Sheet）}\\[18pt]
{\LARGE\bfseries 群落演替模拟：在干旱中生存}\\[8pt]
{\large （Community Succession Simulation: Surviving Drought）}\\[10pt]
\end{center}

\begin{center}
{\large\bfseries 摘\quad 要}
\end{center}

\vspace{6pt}
植物群落的物种丰富度对其抵御干旱的能力具有显著影响。本文从生态位的视角出发，模拟干旱条件下植物群落的演替过程，并分析各类因素的作用。

首先，我们建立了基于生态位理论的植物群落演替模型，用以模拟干旱条件下的群落演替。该模型以生态位宽度的变化作为刻画群落演替的指标，并利用竞争矩阵描述物种之间的相互作用，包括竞争与促进。随后，我们建立生态位宽度模型，综合考虑繁殖、死亡以及不确定性的影响。其中，繁殖取决于干旱程度、生态位宽度、种间关系与环境容量，而死亡取决于干旱程度与生态位宽度。不确定性以高斯白噪声度量。

基于 Beverton--Holt 方程与 Lotka--Volterra 方程，我们构建了生态位宽度微分方程模型。为刻画各种不规则的天气周期，我们采用服从正态分布的随机变量来模拟干旱的程度与持续时间。作为算例，我们模拟了一个三物种群落的演化，并得出结论：群落的总生态位宽度增加，表明其能够长期存续。

为考察物种数量对群落动态的影响，我们利用植物群落演化模型分别模拟了包含 1 至 4 个物种的群落。结果表明，群落要从中获益至少需要三个物种。此外，随着物种数量的增加，抗旱能力先提高而后趋于稳定。

基于竞争矩阵，我们设置了三个不同的群落，并通过模拟讨论了它们的抗旱能力。由此可以得出结论：以促进作用为主导的种间关系能够增强抗旱能力并延长群落的存续时间。此外，我们改变了干旱的发生频率与强度，进而得出结论：物种灭绝更可能发生在严重干旱期间。

污染与生境萎缩会影响植物群落的生态位宽度与竞争关系。因此，我们对生态位宽度与竞争矩阵作了相应调整。结果表明，污染与生境萎缩可能导致某些物种灭绝。

最后，基于各类因素的影响，我们提出了若干确保植物群落长期存续的措施。此外，我们还分析了保护脆弱群落对生态系统整体健康的重要意义。

\vspace{6pt}
\noindent\textbf{关键词：}群落演替；抗旱性；生态位宽度；种群相互作用；Beverton--Holt 方程

\newpage
\pagestyle{fancy}

%%====================== 正文 ======================
\section{引言}

\subsection{背景}

不同物种的植物对干旱的敏感性与抵抗能力存在差异$^{[1]}$。大量观测表明，植物群落的物种丰富度显著影响其在长期内适应水分匮乏的能力$^{[2]}$。含有较多物种的群落在其后续世代中往往表现出更强的干旱胁迫抗性，而物种较少的群落表现出较弱的抗性。因此，分析植物群落的干旱适应性与物种数量之间的关联，对其在较长时期内的存续至关重要。

\subsection{问题重述}

\begin{itemize}
\item 建立一个模型，用以预测植物群落在各种不规则天气周期下的演化，并考虑物种之间的相互作用。
\item 确定群落要从中获益所需的最小物种数，以及物种数量增加对群落的影响。
\item 分析物种类型对群落演化的影响。
\item 讨论干旱发生频率与持续宽度的增大或减小所带来的影响。
\item 分析污染与生境缩减等其他因素对模型的影响。
\item 依据模型确定应当采取何种措施以确保植物群落的长期存续，以及这些措施对更大环境的影响。
\end{itemize}

\subsection{我们的工作}

概括全文，我们做了以下工作：

\begin{itemize}
\item 建立了一个考虑不确定天气周期的生态位模型，用以模拟植物群落演化。该模型基于种间竞争刻画物种的相互作用与演替过程，并建立竞争矩阵来描述群落内部的种群相互作用。
\item 利用该模型确定了群落从物种数量增加中获益所需的最小物种数。模型考虑不确定干旱条件下各种群的生态位宽度，并使用基于 Beverton--Holt 方程与 Lotka--Volterra 方程的微分方程。
\item 分析了不同物种类型对群落演化的影响。模型表明，不同类型的种群具有不同的相互作用，互利共生可能带来更强的抗旱能力。然而，在严重且持久的干旱中，竞争占优可能导致种群灭绝，从而降低物种丰富度与群落稳定性。
\item 利用该模型研究了各种干旱周期对群落演化的影响。分析表明，频繁且持续时间较长的干旱可能对种群产生负面影响，而频率较低的干旱可能使种群更适应干旱环境。
\item 讨论了污染与生境缩减等其他因素对群落演化的影响。模型设定污染系数与生境缩减系数以影响生态位宽度与物种相互作用，从而影响群落演替。污染与生境缩减会加剧竞争，使竞争力更强的种群占据资源，但可能导致物种灭绝并影响群落稳定性。
\item 提出了确保植物群落长期存续的措施及其对更大环境的影响。增加群落中的物种数量可能提高干旱适应性，但必须与避免竞争、确保互利共生以及减少环境污染和生境缩减相平衡，以保证群落的稳定性和长期存续。
\end{itemize}

\section{符号与假设}

\subsection{符号}

\begin{itemize}
\item $C$：物种的竞争矩阵
\item $L$：物种的生态位宽度
\item $\Gamma$：物种对干旱的敏感性
\item $l_i(t)$：第 $i$ 个物种在时刻 $t$ 的生态位宽度
\item $b_i(t)$：第 $i$ 个物种因其繁殖而增加的生态位宽度
\item $d_i(t)$：第 $i$ 个物种因死亡而减少的生态位宽度比例
\item $\sigma_i(t)$：用于刻画不可预测性的高斯白噪声
\item $K(t)$：群落于时刻 $t$ 的生态位宽度
\item $g_i$：第 $i$ 个物种在无竞争条件下的增长率
\item $d_r$：描述干旱程度的干旱系数
\item $f(d_r)$：影响死亡率的干旱响应函数
\item $t_{INR}$：距下一次干旱的时间间隔
\item $t_d$：干旱的持续时间
\end{itemize}

\subsection{假设}

为简化问题并便于我们模拟现实情形，我们作出以下基本假设，每条假设均给出了适当的合理性说明。

\begin{itemize}
\item \textbf{假设 1：}植物群落中的物种数量不会增加。\\
理由：我们假设系统是封闭的，没有新物种被引入，也不会有新物种能够定殖于该植物群落。
\item \textbf{假设 2：}物种不发生突变，但种群及其所控制的资源随时间变化。\\
理由：本研究的所涉时间尺度相对较短，可以认为可能导致生态变化的遗传改变与突变可忽略不计。
\item \textbf{假设 3：}群落内物种之间的竞争关系或互利关系保持不变。\\
理由：群落内物种之间的相互作用十分复杂，并可能因多种因素随时间变化。此处作此处理是为了简化模型，并聚焦于干旱对群落结构的影响。
\item \textbf{假设 4：}群落中每个物种的干旱敏感系数是固定的，不随时间和空间变化。\\
理由：为使模型简单易解，各物种的生态位参数常被视为不变量，这有助于我们更好地理解物种之间的相互作用以及干旱事件对群落结构的影响。
\end{itemize}

\section{植物群落演化模型}

生态位是指物种在群落中相对于相关物种的功能关系与角色，并考虑时间与空间。生态位宽度是生物所利用资源多样性的指标$^{[3]}$。物种的生态位宽度越宽，其专化程度越低，对环境的适应性越强。

\subsection{模型概述}

生态位模型基于种间相互作用描述植物群落的演化$^{[4]}$。依据生态位模型，每个物种占据一个特定的生态位，具有特定的资源利用策略与环境适应性。该模型通过资源竞争与合作来考虑物种的相互作用。

生态位模型模拟植物群落演化的基本步骤如下。

基于生态位模型的群落演化过程如下图所示。

\subsection{物种的初始条件与生态位特征}

在群落演替的初始时刻，共有 $m$ 个物种，环境容量为 $K$。环境容量是指给定环境所能容纳的最大种群数量。

干旱敏感系数是常用于表征植物物种对干旱敏感程度的指标。该系数通常取 $0$ 与 $1$ 之间的数值，数值越大表示对干旱越敏感。

群落的干旱敏感性可由以下向量表示：
\begin{equation}
\Gamma=[\gamma_1,\gamma_2,\cdots,\gamma_m]
\end{equation}
其中 $\gamma_i$ 表示第 $i$ 个物种的干旱敏感系数。

在生态位模型中，竞争矩阵是描述不同物种之间竞争与互利关系的矩阵$^{[5]}$，其表达式如下：
\begin{equation}
C=\begin{bmatrix}
c_{11} & c_{12} & \cdots & c_{1m}\\
c_{21} & c_{22} & \cdots & c_{2m}\\
\vdots & \vdots & \ddots & \vdots\\
c_{m1} & c_{m2} & \cdots & c_{mm}
\end{bmatrix}
\end{equation}
其中 $c_{ij}$ 表示第 $i$ 个物种与第 $j$ 个物种之间的竞争系数。当 $c_{ij}=0$ 时，两个物种之间不存在竞争；$c_{ij}>0$ 表示第 $i$ 个物种损害第 $j$ 个物种的生存与繁殖，即竞争；相反，$c_{ij}<0$ 表示第 $i$ 个物种对第 $j$ 个物种的存在与繁殖具有正向影响，即共生或互利。竞争系数的绝对值越大，竞争或互利关系越强。

竞争矩阵通常需要依据实验数据、生态调查数据或其他相关数据与知识来估计。考虑到物种之间的竞争是双向且相互影响的，竞争矩阵通常是对称矩阵，即 $c_{ij}=c_{ji}$。

在生态学中，生态位是指物种在其生态系统中的功能与作用，即物种在生态系统中的位置、其利用资源的方式，以及物种与其周围环境之间的相互作用。

生态位宽度是指生物所能利用的各种资源的总和，是衡量生物所利用资源多样性的指标。设群落的生态位宽度为：
\begin{equation}
L=[l_1,l_2,\cdots,l_m]
\end{equation}
其中 $l_i$ 表示第 $i$ 个物种所利用资源的总和。$L$ 值越大，物种所能利用的资源越多。

考虑到群落的扩张，各物种生态位之和 $K$ 可在特定范围内变化。
\begin{equation}
K(t)=\sum_{i=1}^{m}l_i(t)
\end{equation}

\subsection{植物群落的演化规则}

生态位宽度受种群个体死亡数与繁殖数的影响。为更准确地模拟气候不稳定与干旱对植物群落演化的影响，在生态位模型中引入随机扰动项。生态位模型如下式所示：
\begin{equation}
l_i(t+1)=l_i(t)+b_i(t)-d_i(t)+\sigma_i(t)
\end{equation}
其中
\begin{itemize}
\item $l_i(t)$ 为第 $i$ 个物种在时刻 $t$ 的生态位宽度；
\item $b_i(t)$ 为第 $i$ 个物种因其繁殖而增加的生态位宽度；
\item $d_i(t)$ 为第 $i$ 个物种因死亡而减少的生态位宽度比例；
\item $\sigma_i(t)$ 为均值为 $0$、方差为 $\sigma^2$ 的高斯白噪声。
\end{itemize}

值得注意的是，在实际应用中，我们需要对随机扰动项作适当调整，以确保其不会对模拟结果产生过大的影响。此外，模型需要进行敏感性分析与验证，以保证其可靠性与可重复性。

\subsubsection{物种繁殖对生态位宽度的影响}

本文中，各物种因繁殖而增加的生态位宽度受当前生态位宽度、环境容量与种间竞争的影响。物种当前的生态位宽度或环境容量越大，其繁殖越多，生态位宽度增加越多。相反，物种之间的竞争越强，其繁殖越少，生态位宽度增加越少。因此，基于 Beverton--Holt 方程$^{[6]}$与 Lotka--Volterra 方程$^{[7][8]}$，得到 $b_i(t)$ 的表达式如下：
\begin{equation}
b_i(t)=g_i l_i(t)\left(1-\frac{l_i(t)+\sum_{j=1,\,j\neq i}^{m}c_{ij}l_j(t)}{K(t)}\right)
\end{equation}
其中 $g_i$ 表示第 $i$ 个物种的非竞争增长率，其仅受环境条件影响。环境条件对该物种越适宜，$g_i$ 越大。

在干旱环境条件下，植物的水资源减少，其增长率相应降低。物种对干旱越敏感，其在干旱条件下的增长率越低。干旱敏感系数并非静止不变，而是随干旱条件变化。物种的非竞争增长率为
\begin{equation}
g_i=g_{\max}e^{-d_r\times\gamma}
\end{equation}
其中 $d_r$ 为干旱系数，用于衡量干旱程度，$\gamma$ 为第 $i$ 个物种的干旱敏感系数，$g_{\max}$ 为其最大增长率。

\subsubsection{物种死亡对生态位宽度的影响}

由死亡导致的生态位宽度减少量与物种的生态位宽度成正比。考虑到干旱期间植物的死亡率升高，用干旱响应函数 $f(d_r)$ 表示物种对干旱程度的响应。该函数可以模拟干旱对植物群落的不利影响。因此，$d_i(t)$ 的表达式如下：
\begin{equation}
d_i(t)=f(d_r)\beta_i l_i(t)
\end{equation}
其中 $f(d_r)$ 描述由干旱导致的物种死亡率，$\beta_i$ 为比例系数。

干旱程度越大，植物物种的死亡率越高。当干旱程度较小时，环境适宜植物生长，此时物种死亡率处于较低水平且变化不大；然而，当干旱严重时，环境恶劣，因此植物的死亡率将处于较高水平。因此，$f(d_r)$ 可基于 Logistic 函数确定。
\begin{equation}
f(x)=\frac{1}{1+e^{-x}}
\end{equation}
令 $x=12(d_r-0.5)$，即把 $d_r$ 由 $[0,1]$ 映射到 $[-6,6]$。
\begin{equation}
f(d_r)=\frac{1}{1+e^{-12(d_r-0.5)}}
\end{equation}
函数 $f(d_r)$ 的图像如下所示。

将式 (6) 与式 (8) 代入式 (5)，得到：
\begin{equation}
l_i(t+1)=l_i(t)+g_i l_i(t)\left(1-\frac{l_i(t)+\sum_{j=1,\,j\neq i}^{m}c_{ij}l_j(t)}{K(t)}\right)-f(d_r)\beta_i l_i(t)+\sigma_i(t)
\end{equation}
进而可得到生态位宽度的微分形式。
\begin{equation}
\frac{\mathrm{d}l_i(t)}{\mathrm{d}t}=g_i l_i(t)\left(1-\frac{l_i(t)+\sum_{j=1,\,j\neq i}^{m}c_{ij}l_j(t)}{K(t)}\right)-f(d_r)\beta_i l_i(t)+\sigma_i(t)
\end{equation}

\subsection{不规则天气周期}

考虑到气候环境的波动、干旱周期与干旱程度的不确定性，我们通过引入随机变量来模拟这些不确定性。气候变化被视为一个随机过程。气候环境波动的模拟基于随机数生成器所产生的随机数。此外，干旱周期与干旱程度服从一定的概率分布。

本文采用指数分布对干旱事件发生的概率建模。指数分布是泊松点过程中事件之间时间间隔的概率分布，即事件以恒定的平均速率连续且独立发生的随机过程。

$\lambda$ 是指数分布的参数，反映干旱事件的发生频率。当 $\lambda$ 较大时，干旱事件的发生频率较高；反之，当 $\lambda$ 较小时，干旱事件的发生频率与程度较低。因此，干旱发生的频率通过 $\lambda$ 的取值来调节。
\begin{equation}
P(t_{INR})=1-e^{-\lambda t_{INR}}
\end{equation}
其中 $P(t_{INR})$ 表示在时间间隔 $t_{INR}$ 内发生干旱事件的概率。

每一次干旱事件的程度与持续时间各不相同。为简化问题，我们认为干旱的程度与持续时间满足相互独立的一维正态分布。
\begin{equation}
d_r\sim N\left(\mu_d,\sigma_d^2\right)\qquad t_d\sim N\left(\mu_t,\sigma_t^2\right)
\end{equation}
其中 $d_r$ 为干旱系数，表示干旱的程度，$t_d$ 为干旱的持续时间。

通过调整干旱发生的概率、强度与持续时间，模型能够模拟变化的气候条件对植物群落演替的影响。当不发生干旱时，$d_r$ 取给定的默认值。气候变化过程如下图所示。
\subsection{植物群落演化的模拟}

\subsubsection{植物物种与初始参数}

本文中，所模拟的群落演化总时长为 100 年，时间步长为 1 年。此外，为简化问题，群落初始含有三个物种。对每个物种而言，其干旱敏感系数、初始生态位宽度与最大增长系数均由表 1 给出。具体数据可通过野外调查或依据文献获得，以下数据仅用于展示我们的模型是如何应用的。

三个物种的相关参数为：物种 1 的干旱敏感系数 $\gamma_1=0.5$、最大增长率 $g_{\max}=0.23$、比例系数 $\beta_1=0.1$、初始生态位宽度 $l_1=2$；物种 2 的 $\gamma_2=0.7$、$g_{\max}=0.15$、$\beta_2=0.1$、$l_2=2$；物种 3 的 $\gamma_3=0.2$、$g_{\max}=0.33$、$\beta_3=0.1$、$l_3=2$。

生态系统的参数为：
\begin{equation}
\begin{cases}
C=\begin{bmatrix}
0 & -0.2 & 0.4\\
-0.2 & 0 & 0.6\\
0.4 & 0.6 & 0
\end{bmatrix}\\[1.0em]
K=10\\[0.3em]
d_r\sim N(1.5,\,0.3^2)\\[0.3em]
t_d\sim N(15,\,5^2)\\[0.3em]
P(t;\lambda=0.04)=1-e^{-\lambda t}
\end{cases}
\end{equation}

\subsubsection{不规则天气的模拟}

依据第 3.4 节关于天气模拟的规则，可得到未来 100 个月的天气。不规则天气模拟的算法如下所示。

\noindent\textbf{算法 1\quad 不规则天气模拟}

\noindent\textbf{输入：}最大迭代次数 $n$

\noindent\textbf{输出：}特定时间点处的干旱系数 $d_r$

\begin{enumerate}
\item $P(t_{INR})\leftarrow 1-\exp(-\lambda t_{INR})$，$\lambda=0.04$；\quad（干旱发生的概率）
\item $d_{r0}\leftarrow 0.2$；
\item $t_{INR}\leftarrow 0$；
\item $i\leftarrow 0$；
\item \textbf{while} $i\le n$ \textbf{do}
  \begin{enumerate}
  \item $r\leftarrow \mathrm{rand}()$；
  \item $p\leftarrow \mathrm{Exponential}(\lambda,\,t_{INR})$；
  \item \textbf{if} $r\le p$ \textbf{then}
     \begin{enumerate}
     \item $day_n\leftarrow \mathrm{round}(2*\mathrm{randn}()+10)$；\quad（$day_n$ 为干旱的持续时间）
     \item \textbf{for} $j\leftarrow i:\,i+day_n-1$ \textbf{do}
        \begin{enumerate}
        \item $d_{rj}=0.15*\mathrm{randn}()+0.6$；
        \end{enumerate}
     \item \textbf{end for}
     \item $i\leftarrow i+day_n$；\quad（更新时间区间端点）
     \item $t_{INR}\leftarrow 0$；\quad（把 $t_{INR}$ 重置为 0）
     \end{enumerate}
  \item \textbf{else}
     \begin{enumerate}
     \item $d_{ri}\leftarrow d_{r0}$；
     \item $t_{INR}\leftarrow t_{INR}+1$；\quad（更新距上次干旱的时间间隔）
     \item $i\leftarrow i+1$；\quad（更新时间点）
     \end{enumerate}
  \item \textbf{end if}
  \end{enumerate}
\item \textbf{end while}
\end{enumerate}

利用不规则天气模拟算法，可以获得未来 100 年的天气，其结果如下图所示。

如图 9 所示，起初气候处于正常状态，环境适宜。随着时间的推移，植物群落将面临不规则的干旱，其持续时间与程度各不相同。

\subsubsection{物种演化的模拟结果}

依据式 (12)，可得到各物种生态位宽度的差分方程。由此，我们通过求解微分方程得到各物种在 100 年内的生态位宽度。我们有
\begin{equation}
\begin{cases}
\dfrac{\mathrm{d}l_i(t)}{\mathrm{d}t}=g_i l_i(t)\left(1-\dfrac{l_i(t)+\sum_{j=1,\,j\neq i}^{m}c_{ij}l_j(t)}{K(t)}\right)-f(d_r)\beta_i l_i(t)+\sigma_i(t)\\[1.2em]
l_i(0)=l_{i0}
\end{cases}
\end{equation}
其中 $\sigma_i(t)$ 服从 $N(0,\,0.5^2)$。

三个物种在 100 年内的生态位宽度如下图所示。

由图 11 可见，三个物种的生态位宽度在初始时刻均等于 2。在不规则天气作用下，物种 1 与物种 2 的生态位宽度在波动中增大，其中物种一的增加最快。总体而言，物种三的生态位宽度先增大、后减小，最终趋于平稳。干旱对几乎所有物种都是不利的。

在不规则天气条件下，各物种生态位宽度所占的比例也随时间变化。图 12 给出了第 1、20、40、60、80 与 100 年各物种生态位宽度的占比。

如图 12 所示，在第一年，我们假设三个物种的生态位宽度占比相等。随着时间推进，物种 1 生态位宽度的占比增加。到第 60 年，各物种生态位宽度的占比趋于稳定。

植物群落的演化主要受干旱影响。图 13 展示了干旱发生对群落总生态位宽度的影响。

如图 13 所示，可以看到在未发生干旱时，群落的总生态位宽度呈增长趋势。这表明此时环境适宜植物生长，种群在此时期的繁殖较多。然而，当干旱发生时，总生态位宽度呈下降趋势。此外，干旱越严重，即 $d_r$ 越大，总生态位宽度下降得越快。这表明干旱抑制种群的繁殖，且抑制效应随干旱程度的增加而增强。

\section{植物群落演化模型的应用}

\subsection{物种数量对群落的影响}

在自然环境中，物种彼此之间既存在竞争关系，也存在互利关系。当环境更严酷时（例如干旱），物种之间的互利关系会为了生存而增强，从而使群落的整体抗旱能力增强。为考察物种数量对群落的影响，我们利用第 3 节的植物群落演化模型，分别模拟了包含 1 至 4 个物种的群落演化。

\subsubsection{群落获益所需的最小物种数}

依据第 3 节的植物群落演化模型，可以模拟包含 1、2、3 与 4 个物种的群落的演化，其结果如下图所示。

由图 14 可以看到，当群落中只有一个物种时，该物种将在不规则干旱条件下缓慢灭绝。当有两个物种时，各物种的生态位宽度随干旱的发生而剧烈波动。这表明面对不规则干旱时，种群无法长期稳定存活。当有三个物种时，各物种的生态位宽度在波动中增大。这表明该群落能够通过种间互利提高抗旱能力，从而使群落得以长期存活。四个物种的情形与三个物种的情形相似。因此可以得出结论：群落从中获益所需的最小物种数为三。

\subsubsection{物种数量增加对群落的影响}

依据 4.1.1 节中四个群落演化的模拟结果，我们可以得出如下结论：随着物种数量的增长，群落的生态位宽度越大，其抗旱能力越强。当物种数量增加到某一数值时，其抗旱能力趋于稳定，甚至可能下降。

单物种群落往往比多物种群落的抗旱能力更弱。多物种群落往往具有更高的资源利用效率、促进效应与物种多样性效应，因此能够更好地适应干旱环境，表现出更强的抗旱能力。其原理如下：

\begin{itemize}
\item 生态位分化：指不同物种之间在资源利用上的差异，从而避免植物之间的竞争。它使群落能够充分利用有限的资源，并增强其抵御干旱的能力。
\item 促进效应：多个物种之间可能存在互利共生，从而提高水分与养分的利用效率，增强植物对干旱的抵抗能力。
\item 物种多样性效应：多样性是衡量生态系统中物种数量与多样程度的指标。多样性的提高可以增强系统的稳定性，减少其受到的外部干扰，从而增强群落的抗旱能力。
\end{itemize}

然而，从生态位模型的角度看，物种过多的群落其抗旱能力未必更强。物种过多会导致生态位重叠与更激烈的种间竞争，还会导致群落稳定性下降，从而使群落的抗旱能力减弱。

\subsection{物种类型对群落演化的影响}

不同植物群落中的物种类型各不相同，因此物种之间的相互作用也各不相同。有些群落中物种之间的竞争更为激烈，而另一些群落中种间互利占主导地位。本文基于第 3 节的物种竞争关系，通过改变竞争矩阵设计了竞争型群落与互利型群落。

\begin{itemize}
\item 一般植物群落\\
该植物群落的物种之间既存在以竞争为主导的关系，也存在以促进为主导的关系，由第 3 节的竞争矩阵表示。
\begin{equation}
C=\begin{bmatrix}
0 & -0.2 & 0.4\\
-0.2 & 0 & 0.6\\
0.4 & 0.6 & 0
\end{bmatrix}
\end{equation}
\item 竞争型植物群落\\
竞争型植物群落以物种之间的竞争关系为主导。其竞争矩阵如下：
\begin{equation}
C=\begin{bmatrix}
0 & 1.0 & 0.8\\
1.0 & 0 & 0.9\\
0.8 & 0.9 & 0
\end{bmatrix}
\end{equation}
\item 互利型植物群落\\
互利型植物群落以物种之间的互利关系为主导。其竞争矩阵如下：
\begin{equation}
C=\begin{bmatrix}
0 & -0.3 & -0.2\\
-0.3 & 0 & -0.25\\
-0.2 & -0.25 & 0
\end{bmatrix}
\end{equation}
\end{itemize}

三个群落的演化如下图所示。可以看到，以互利为主导的群落将具有更大的生态位宽度与更多的物种繁殖。这表明其抗旱能力更强。相反，当群落中物种之间存在激烈竞争时，其生态位宽度更小，种群繁殖更少，表明其抗旱能力更弱。

\subsection{干旱严重程度的影响}

干旱程度与持续时间的增加会对植物群落产生负面影响。干旱的严重程度可通过调整植物群落演化模型中的干旱系数 $d_r$ 与干旱持续时间 $t_d$ 来模拟。当干旱变得更为严重时，$d_r$ 与 $t_d$ 服从以下正态分布：
\begin{equation}
d_r\sim N\left(0.8,\,0.1^2\right)\qquad t_d\sim N\left(12,\,3^2\right)
\end{equation}
当干旱较为温和时，$d_r$ 与 $t_d$ 服从以下正态分布：
\begin{equation}
d_r\sim N\left(0.45,\,0.15^2\right)\qquad t_d\sim N\left(8,\,1.5^2\right)
\end{equation}
基于植物群落演化模型，三种干旱条件下的群落演替过程如下图所示。

由图 16 可以看到，随着干旱加剧，群落的生态位宽度减小，这意味着有更多物种死亡。因此，群落能够存续的时间更短。

\subsection{污染与生境缩减的影响}

污染与生境萎缩会降低生态位宽度，限制生物可获得的资源数量。资源的这一减少会加剧竞争强度，使竞争力更强的物种占据主导地位，并对其他物种施加额外的压力，导致其衰退或灭绝。

在植物群落演化模型中，灭绝的物种将从竞争矩阵中被移除。这可能导致竞争关系发生变化，进而影响其他物种的竞争能力以及群落演替过程。

污染与生境萎缩会影响植物群落的生态位宽度与竞争关系。因此，可对生态位宽度与竞争矩阵作如下调整。
\begin{equation}
L'=L(1-P_r)H\qquad C'=C(1-P_r)H
\end{equation}
其中 $L'$ 表示调整后的生态位宽度，$C'$ 表示调整后的竞争矩阵，$P_r$ 表示植物群落的污染程度，$H$ 表示生境缩减系数。

依据植物群落演化模型，可以模拟面临污染与生境萎缩的群落的演化。其结果如下图所示。

如图 17 所示，当污染较轻、生境缩减较小时，群落的生态位宽度相对减小。此时植物群落仍然能够长期稳定存活。当污染严重、生境缩减较大时，生态位宽度接近于 0，即此时群落中的许多物种将会灭绝。

\subsection{使植物群落存续更久的措施}

对生态位模型的分析表明，多样性较高的植物群落在干旱环境中具有更强的抗旱能力。这是由于不同物种的生态位互补，减少了竞争，使其能够更好地适应干旱带来的挑战。为确保植物群落的长期存续，可采取以下措施：

\begin{enumerate}
\item 保护自然环境，维护生态系统的完整性与稳定性，减少人类活动对自然环境的破坏。尤其应当保护植被与水源，防止过度采伐与过度放牧、化肥与农药的过度使用以及水资源的过度开发等。
\item 鼓励植物多样性，维持群落的多样性。通过种植不同种类的植物，可以增加植物群落的多样性，提高植物的适应性与抗性，减少植物之间的竞争，并维持群落的稳定性。
\item 推行节水措施。水是维持生命的重要资源。人类可以采取合理利用与节约水资源、减少对天然水源的污染与破坏等节水措施，为植物提供更好的生长环境。
\item 培育抗旱品种。可通过育种或基因编辑技术培育抗旱植物品种，以提高植物的抗旱能力，使其更好地适应干旱环境。
\end{enumerate}

综上所述，生态位模型表明，生物群落中的每个物种都具有特定的生态位，并通过竞争与互利彼此相互作用。保护脆弱群落能够维持生态系统的稳定性与生态多样性。此外，保护脆弱群落对人类自身利益也至关重要，因为生态系统中的动植物为人类的需求提供资源，而对生态系统的破坏会影响人类的生产与生活。因此，保护脆弱群落对于维持生态系统稳定、保护生态多样性以及保障人类自身利益都具有重要意义。

\section{敏感性分析}

\subsection{对不可预测因素的敏感性}

在建立种群生态位宽度的微分方程时，我们加入了白噪声，以模拟不可预测因素对群落演替的影响。白噪声的方差需要处在一定范围之内。若方差过小，则无法模拟不可预测因素的影响；若方差过大，则会影响群落演替过程的正确性。

白噪声的方差对群落的整体生态位宽度影响较小。这是因为我们对每个物种的白噪声是独立计算的，当物种较多或模拟次数足够时，群落的整体生态位宽度变化很小。然而，某些偶然因素可能对群落中的所有物种产生负面影响，而不是独立地作用于每个物种。因此，在考虑不可预测干扰时，需要更复杂的模型才能获得更好的结果。

\subsection{对干旱系数的敏感性}

干旱系数用于描述物种对干旱的敏感性，我们通过调整干旱系数来观察具有不同干旱适应性的物种与群落的表现。

从图中可以看出，物种对干旱越敏感，其生态位宽度越小，整个群落的生态位宽度也越小。

\section{优缺点}

\subsection{优点}

\begin{itemize}
\item 生态位模型能够全面刻画不同物种之间的生态关系，包括竞争与合作，并能够模拟植物群落随时间的演化。
\item 该模型具有灵活性，可以适应不同的环境条件，例如干旱事件或生境丧失，从而能够评估这些干扰对群落的潜在影响。
\item 生态位模型能够纳入关于植物物种相互作用的经验观测，因此是检验生态学理论、预测群落对环境变化响应的有用工具。
\end{itemize}

\subsection{缺点}

\begin{itemize}
\item 该模型可能对群落的初始条件较为敏感，这些条件的微小变化可能导致群落演化随时间出现显著差异。
\item 该模型假设物种相互作用仅基于对资源的竞争与合作，这可能无法刻画群落中的所有生态相互作用，例如捕食或寄生。
\item 竞争矩阵参数所需满足的条件可能与自然界中的实际情况不一致，例如竞争矩阵未必是对称矩阵。
\end{itemize}

\section{结论}

\begin{itemize}
\item 由于资源利用效率更高以及物种多样性效应，多物种群落通常比单物种群落表现出更强的抗旱能力。
\item 当物种数量增加到一定程度时，其抗旱能力趋于稳定甚至下降。这是因为物种过多会导致生态位重叠，从而加剧种间竞争。
\item 以互利关系为主导的群落具有更强的干旱适应能力。
\item 当干旱发生得更频繁、持续时间更长时，物种可能无法迅速适应变化的环境，导致某些物种的种群规模减小，从而破坏生态平衡。相反，若干旱发生频率较低，物种则更有可能适应干旱环境，因为它们有更多时间来适应干旱条件。
\item 生境污染与生境缩减都会导致生态位宽度减小，加剧竞争程度。竞争力强的物种可能占据更多资源，对其他物种施加更大的竞争压力，导致其他物种减少或灭绝。
\end{itemize}

\section*{参考文献}

\begin{itemize}
\item[{[1]}] DT Rosenow et al. ``Drought tolerant sorghum and cotton germplasm''. In: \textit{Agricultural Water Management} 7.1--3 (1983), pp. 207--222.
\item[{[2]}] J Sardans and Josep Peñuelas. ``Plant-soil interactions in Mediterranean forest and shrublands: impacts of climatic change''. In: \textit{Plant and soil} 365 (2013), pp. 1--33.
\item[{[3]}] Mustaqeem Ahmad et al. ``Niche width analyses facilitate identification of high-risk endemic species at high altitudes in western Himalayas''. In: \textit{Ecological Indicators} 126 (2021), p. 107653.
\item[{[4]}] Jitka Polechová and David Storch. ``Ecological niche''. In: \textit{Encyclopedia of ecology} 2 (2008), pp. 1088--1097.
\item[{[5]}] Andrew D Letten, Po-Ju Ke, and Tadashi Fukami. ``Linking modern coexistence theory and contemporary niche theory''. In: \textit{Ecological Monographs} 87.2 (2017), pp. 161--177.
\item[{[6]}] M De la Sen. ``The generalized Beverton--Holt equation and the control of populations''. In: \textit{Applied Mathematical Modelling} 32.11 (2008), pp. 2312--2328.
\item[{[7]}] S Das and PK Gupta. ``A mathematical model on fractional Lotka--Volterra equations''. In: \textit{Journal of Theoretical Biology} 277.1 (2011), pp. 1--6.
\item[{[8]}] Martin Bohner and Rotchana Chieochan. ``The Beverton--Holt $q$-difference equation''. In: \textit{Journal of Biological Dynamics} 7.1 (2013), pp. 86--95.
\end{itemize}

\end{document}
